% TOP_PROBLEM: {"id": 30007254, "title": "Zilber-Pink for abelian varieties over Qbar", "category_id": 1, "category": {"id": 1, "name": "number_theory", "display_name": "Number Theory", "description": "Properties of integers, prime numbers, Diophantine equations.", "slug": "number-theory", "order_index": 1, "created_at": "2026-07-31T15:26:25.670Z"}, "status": "open"}
% FIRST_POSTED: 2026-10-06
% ATTEMPT_STATUS: solved
% OPENAI_CLAIM: {"schema_version":1,"source_repo":"https://github.com/openai/math","source_commit":"adc7f1241b42e322a6451854ab7e4b4c146bf78a","release_date":"2026-10-06","match":"full","resolution":"proved","independently_reviewed":false,"second_pass":"agreed","families":[{"family":"016","title":"Zilber–Pink in abelian varieties and the Siegel threefold","subject":"Number theory","manuscripts":[{"title":"The abelian Zilber–Pink conjecture","pdf_path":"preprints/The-Abelian-Zilber-Pink-Conjecture-September-24-2026/paper.pdf","readme_path":"preprints/The-Abelian-Zilber-Pink-Conjecture-September-24-2026/README.md","date":"2026-09-24","theorem_ref":"thm:main (sections/01-introduction.tex:43)","statement_sources":[{"path":"preprints/The-Abelian-Zilber-Pink-Conjecture-September-24-2026/build/sections/01-introduction.tex","line":43,"label":"thm:main"}]}],"lean":null,"reasoning_trace":null}]}
\documentclass[11pt]{article}
\usepackage[margin=25mm]{geometry}
\usepackage{fontspec}
\setmainfont{Latin Modern Roman}
\usepackage{amsmath,amssymb,mathtools}
\usepackage{unicode-math}
\setmathfont{Latin Modern Math}
\usepackage{enumitem}
\usepackage{xurl,hyperref,bookmark}
\hypersetup{colorlinks=true,urlcolor=blue}
\setcounter{secnumdepth}{0}
\setlength{\parindent}{0pt}
\setlength{\parskip}{5pt}
\setlength{\emergencystretch}{3em}
\begin{document}
\section*{Zilber-Pink for abelian varieties over Qbar}

\subsection{OpenAI's claim}
OpenAI presents the following claims relevant to this question, with manuscript numbers following paper order in Sources. In manuscript 1 (family 016; 2026-09-24), the asserted result is that subvarieties of abelian varieties over algebraic numbers have only finitely many maximal atypical intersections, measured inside their smallest containing torsion coset.

\subsection{Mathematical statement of the claimed result}
\paragraph{Manuscript 1.}
For a number field $K$, abelian variety $A/K$, and irreducible $X\subseteq A_{\bar K}$, let $S$ be the smallest special subvariety containing $X$; special subvarieties are torsion translates of abelian subvarieties. An irreducible component $Y$ of $X\cap Z$, $Z\subseteq S$ special, is atypical when
\[
 \dim Y>\dim X+\dim Z-\dim S.
\]
The claim is finiteness of maximal atypical subvarieties of $X$, equivalently existence of finitely many proper atypical $Y_j\subsetneq X$ containing every atypical subvariety. No dimension bound on $A$ or $X$ is imposed.

\subsection{Outline of the argument}
For manuscript 1 (family 016; 2026-09-24), a counterexample sequence is decomposed into distinct height scales. Uniform tangent positivity and weighted small sections give a height inequality despite growing homomorphism coefficients and point degrees; passing to suitable quotients rules out the remaining scale configurations.

\subsection{Scope relative to this problem}
The result supplies the conclusion in this problem with its displayed hypotheses and quantifiers. The broader statements and applications in the manuscript are kept separate from the particular question above. The supporting result is manuscript 1 (family 016; 2026-09-24).

\subsection{Formal verification}
No Lean formalization is listed at this snapshot for manuscript 1 (family 016; 2026-09-24).

This site has not executed the Lean code. Described statement coverage is not an independent proof check.

\subsection{Status caveat}
The manuscripts were produced by an internal OpenAI model. OpenAI cautions that results without formalization can contain mistakes. This site has not independently checked the proofs; the label records the scope of OpenAI's claim.

\subsection{Sources}
\begin{itemize}
\item OpenAI, \emph{The abelian Zilber–Pink conjecture}, 2026-09-24; citation key \texttt{OAI:The-Abelian-Zilber-Pink-Conjecture-September-24-2026}. \href{https://github.com/openai/math/blob/adc7f1241b42e322a6451854ab7e4b4c146bf78a/preprints/The-Abelian-Zilber-Pink-Conjecture-September-24-2026/paper.pdf}{PDF (GitHub)}; \href{https://github.com/openai/math/raw/adc7f1241b42e322a6451854ab7e4b4c146bf78a/preprints/The-Abelian-Zilber-Pink-Conjecture-September-24-2026/paper.pdf}{PDF (direct)}; \href{https://github.com/openai/math/blob/adc7f1241b42e322a6451854ab7e4b4c146bf78a/preprints/The-Abelian-Zilber-Pink-Conjecture-September-24-2026/README.md}{Manuscript page and citation}; \href{https://github.com/openai/math/blob/main/preprints/The-Abelian-Zilber-Pink-Conjecture-September-24-2026/paper.pdf}{Latest version}; Statement source: \href{https://github.com/openai/math/blob/adc7f1241b42e322a6451854ab7e4b4c146bf78a/preprints/The-Abelian-Zilber-Pink-Conjecture-September-24-2026/build/sections/01-introduction.tex\#L43}{thm:main}.
\item \href{https://github.com/openai/math/blob/adc7f1241b42e322a6451854ab7e4b4c146bf78a/CONTENTS.md}{OpenAI catalogue entry}.
\end{itemize}
\end{document}
